% Chapter 1: What Is Computation?
% Part I: The Nature of Computation

\chapter{What Is Computation?}
\label{ch:computation}

%======================================================================
% SECTION 1: THE QUESTION
%======================================================================
\section{The Question: What Problem Was Humanity Trying to Solve?}
\label{sec:computation:question}

In 1928, David Hilbert posed the \emph{Entscheidungsproblem} (decision problem): Is there a mechanical procedure that, given any mathematical statement, can decide whether it is true or false?

This was the culminating question of Hilbert's program --- the dream of placing all mathematics on a completely formal, axiomatic foundation where truth becomes mechanically decidable. If such a procedure existed, mathematics would be reduced to calculation. No intuition, no creativity, no genius required --- just a machine following rules. Hilbert famously declared: ``We must know --- we shall know'' (\emph{Wir m\"ussen wissen --- wir werden wissen}).

But the question presupposed an answer to a deeper question: \emph{What does it mean for a procedure to be ``mechanical''?} Before we can ask whether a mechanical procedure exists, we must define what counts as a mechanical procedure. This is the question that birthed computer science.

What is a computation?

What distinguishes a ``mechanical procedure'' that a human computer could follow from a mathematical insight that requires understanding? Can every well-defined mathematical procedure be captured by a single, universal formalism?

To appreciate the depth of this question, consider a simple example. Suppose we want to decide whether a Diophantine equation --- an integer polynomial equation like $x^3 + y^3 = z^3$ --- has a solution in integers. A trained mathematician might stare at it, summon deep number-theoretic insights, and eventually recognize that Fermat's Last Theorem (proved by Wiles in 1994) implies no solution exists. But a mechanical procedure would need to reach the same conclusion by systematically applying predetermined rules, without any understanding of number theory. The question is: can \emph{every} such mathematical decision be reduced to rule-following?

The Entscheidungsproblem was the third pillar of Hilbert's program. The first pillar was \emph{consistency}: that a formal system for mathematics cannot prove a contradiction. The second was \emph{completeness}: that every true mathematical statement can be proved within the system. The third was \emph{decidability}: that there exists an algorithm to determine truth or falsity of any statement. G\"odel's incompleteness theorems (1931) had already shattered the first two pillars. The Entscheidungsproblem was the last remaining hope for Hilbert's vision.

\subsection{Why Humans Wanted Mechanical Reasoning}
\label{sec:computation:why_mechanical}

The dream of mechanical reasoning is as old as logic itself. Aristotle's syllogisms (\emph{Organon}, circa 350 BCE) were presented as rules that guarantee truth if followed correctly --- the first suggestion that reasoning is not a mysterious gift but a \emph{process} that can be followed. Two millennia later, the ambition sharpened. Ramon Llull's \emph{Ars Magna} (1305) built literal ``thinking wheels'' for combining concepts mechanically. Descartes and Pascal mechanized arithmetic --- Pascal's calculator (1642) added and subtracted automatically --- but arithmetic was the easy case; the hard problem was reasoning itself. It was Gottfried Wilhelm Leibniz who took the decisive step. In his \emph{Dissertatio de Arte Combinatoria} (1666), Leibniz proposed a \emph{characteristica universalis}, a universal symbolic language in which every concept is decomposed into atomic primitives, together with a \emph{calculus ratiocinator}, a calculus of reasoning over those symbols. His ambition: any dispute, including a disagreement between two philosophers, could be ``settled as easily as a dispute between two accountants.'' Truth would be \emph{calculated}: \emph{calculemus!} --- ``let us calculate.'' Leibniz even devised a binary notation, glimpsing that symbols themselves could be manipulated mechanically. This is the ambition Hilbert inherited: the reduction of reasoning to rule-following.

\subsection{Why Certainty Mattered}
\label{sec:computation:why_certainty}

Why did this question carry such weight? Because mathematics was, for most of its history, the gold standard of certainty --- the one domain where disagreement could in principle be settled. Euclid's \emph{Elements} stood for two millennia as the model of incontrovertible proof. The nineteenth century shattered that calm. Non-Euclidean geometry (Lobachevsky, Bolyai, Gauss) showed that even the axioms of space could be varied. Worse, in 1901 Bertrand Russell found a paradox at the very foundations of set theory: the set of all sets that do not contain themselves both does and does not contain itself. Frege's \emph{Grundgesetze der Arithmetik} was already in press; Frege appended a note: ``A scientist can hardly meet with anything more undesirable than to have the foundation give way just as the work is finished.'' If the foundations of mathematics harbored a contradiction, then \emph{every} proof --- including the scientific conclusions built upon mathematics --- rested on sand. The \emph{Grundlagenkrise} (crisis of foundations) was not academic pedantry: it was the demand that mathematics deserve its authority. Restoring that certainty by purely formal, finitary means was the aim of Hilbert's program, and the question whether such means could be \emph{mechanized} is the question that birthed computer science.

%======================================================================
% SECTION 2: WHY IT SEEMED IMPOSSIBLE
%======================================================================
\section{Why It Seemed Impossible}
\label{sec:computation:impossible}

In the 1920s and 1930s, defining ``computation'' formally seemed paradoxical. Computation was something \emph{people did} --- inherently tied to understanding, meaning, and intention. The idea that it could be captured by a mathematical object --- a machine manipulating symbols without understanding them --- struck many as a category error.

``A machine cannot think. A machine follows rules. But mathematics requires insight. Therefore no machine can do mathematics.'' The gap between \emph{syntax} (symbol manipulation) and \emph{semantics} (meaning) seemed unbridgeable.

Three competing intuitions blocked progress:

\begin{enumerate}
\item \textbf{The Human Computer Model:} Computation is what a human clerk does with paper and pencil. But human computers make mistakes, get tired, need training. How do you formalize ``a human following instructions'' without formalizing the human? A human brings background knowledge, judgment, and creativity to every calculation. Turing's key insight was to analyze \emph{what aspects of the human computer are essential} and which are incidental --- removing the humanity while keeping the computation.

\item \textbf{The Function Model:} Computation is evaluating a mathematical function. But which functions? Only computable ones --- but that's circular. The class of ``effectively calculable functions'' had no formal definition. Moreover, not every function is computable: consider the function that maps a real number to its integer part (computable) versus one that maps a real number to $1$ if a given Turing machine halts and $0$ otherwise (not computable). The boundary between the computable and the uncomputable was entirely unknown.

\item \textbf{The Algorithm Model:} An algorithm is a finite list of instructions. But what counts as an instruction? ``Multiply two numbers'' --- is that one instruction? What about ``solve this differential equation''? The granularity problem seemed insoluble. If instructions are too coarse, the algorithm hides complexity. If too fine, the description becomes impractically large. What is the right level of atomicity for a ``mechanical step''?
\end{enumerate}

Consider the Euclidean algorithm for GCD: ``Take the larger number, subtract the smaller from it repeatedly until the remainder is less than the smaller.'' Is ``subtract repeatedly'' one instruction or many? If the latter, how do we express iteration without presupposing a mechanism for looping? The problem is that any sufficiently expressive instruction set already contains the computational power we are trying to define.

Experts believed any formal definition would either be too narrow (excluding clearly computational processes) or too broad (including things that aren't). Capturing the intuitive notion exactly seemed impossible --- like trying to formalize ``art'' or ``beauty.''

In 1934, G\"odel rejected Church's proposal that $\lambda$-definability captures effective calculability, calling it ``thoroughly unsatisfactory.'' G\"odel demanded a justification that the proposed formal class could not be extended by any intuitively sound principle. Turing's analysis, two years later, would provide exactly that justification.

%======================================================================
% SECTION 3: HISTORICAL CONTEXT
%======================================================================
\section{Historical Context: What Was Known at the Time}
\label{sec:computation:context}

By 1928, mathematics was in crisis. Russell's paradox (1901) had shattered naive set theory. G\"odel's completeness theorem (1929) showed first-order logic is complete --- but his incompleteness theorems (1931) would soon show that any sufficiently strong formal system contains unprovable truths. Hilbert's program --- to prove mathematics consistent, complete, and decidable using finitary methods --- was the last great hope for mathematical certainty. To see why that program took the form it did --- and why answering it required a definition of \emph{computation} --- it helps to follow the two-millennia thread that led to the Entscheidungsproblem.

\subsection{From Euclid to Leibniz}
\label{sec:computation:euclid_to_leibniz}

The search for a mechanical method of reasoning begins with the axiomatic method itself. Euclid's \emph{Elements} (circa 300 BCE) organized geometry as a small set of axioms and rules of inference, showing that a vast body of truth could be generated from a few starting points. al-Khw\=arizm\=i's algebra (circa 820 CE) gave the world the word \emph{algorithm} --- a named, repeatable procedure. Yet for two millennia the axiomatic method was confined to the \emph{content} of mathematics; the act of reasoning itself was left to human judgment. The seventeenth century changed the ambition. Descartes dreamed of a universal method for acquiring knowledge; Pascal built the first mechanical calculator (1642); and Leibniz --- reading both --- proposed to do for reasoning what Pascal had done for arithmetic: reduce it to symbolic manipulation. His \emph{characteristica universalis} (a universal symbolic language) together with his \emph{calculus ratiocinator} (a calculus of reasoning over those symbols) would, he hoped, settle any dispute ``as easily as a dispute between two accountants.'' \emph{Calculemus!} --- ``let us calculate.'' Leibniz even devised a binary notation, glimpsing that symbols themselves could be manipulated mechanically.

\subsection{The Birth of Mathematical Logic}
\label{sec:computation:birth_of_logic}

Leibniz's dream was not realized in his lifetime; realizing it took two centuries of accumulating precision. The key step was to make \emph{logic itself} a mathematical object. George Boole's \emph{The Laws of Thought} (1854) turned propositions into an algebra of $0$ and $1$ --- logical reasoning as arithmetic. Gottlob Frege's \emph{Begriffsschrift} (1879) introduced quantifiers and a formal predicate logic --- the first system in which the entire structure of a mathematical proof could be written out symbolically. Peano's axioms (1889) showed that the natural numbers could be characterized by finitely many postulates. Russell and Whitehead's \emph{Principia Mathematica} (1910--13) then spent three volumes deriving arithmetic from logic --- a demonstration of how far formalization could go. By the early twentieth century, ``logic'' was no longer a branch of philosophy; it was a mathematical discipline with an exact syntax. The stage was set for the question Hilbert would pose: could the whole of mathematics be captured by such formal rules?

\subsection{Hilbert's Grand Vision}
\label{sec:computation:hilbert_vision}

David Hilbert, the leading mathematician of his generation, saw in this formalization an opportunity to secure mathematics once and for all. His 1900 list of 23 problems opened with the consistency of arithmetic. Over the following decades, Hilbert developed his \emph{program}, built on two convictions. First, \emph{formalization}: every mathematical statement and proof can be written in a precise formal language with explicit rules of inference. Second, \emph{finitism}: the metamathematical investigation of such a system --- the proof that it cannot produce a contradiction --- must itself be conducted by finitary, combinatorial reasoning, the kind that does not presuppose the mathematics being justified. If consistency could be proved this way, mathematics would be certified from the outside, without circularity. Hilbert's confidence was total; he declared that ``in mathematics there is no \emph{ignorabimus}'' (``there is no ``we shall not know''), and closed his 1930 radio address with ``\emph{Wir m\"ussen wissen --- wir werden wissen}'' (``We must know --- we shall know'').

\subsection{Logic Versus Intuition}
\label{sec:computation:logic_vs_intuition}

Hilbert's program was not universally accepted; it was the flashpoint of the foundational crisis. Its leading opponent was L.~E.~J. Brouwer, founder of \emph{intuitionism}. Brouwer denied that a mathematical object exists unless it can be constructed --- and therefore rejected the law of the excluded middle (every statement is either true or false) in its unrestricted form. For Brouwer, Hilbert's formal proofs about meaningless symbol strings were not mathematics at all; mathematics is the free creation of the mind, prior to language. Hilbert, by contrast, was content to treat formal symbols as meaningless tokens: ``we can operate with formulas in the same way that a chess player operates with the pieces.'' The dispute over the excluded middle was not idle: a classical proof of the form ``either $P$ or not $P$'' may establish existence without exhibiting a construction, and for an intuitionist that is no proof at all. The controversy matters for computation because it forced a question Hilbert could defer: what counts as a legitimate proof \emph{step}? A finitary, mechanical step --- or something requiring insight? The notion of an ``effective procedure'' stayed unsettled precisely until Turing made it precise.

\subsection{The Entscheidungsproblem}
\label{sec:computation:entscheidungsproblem}

The decision problem emerged explicitly in Hilbert and Ackermann's \emph{Grundz\"uge der theoretischen Logik} (1928), which asked whether there is an effective procedure that decides, for any formula of first-order logic, whether it is valid. Once G\"odel's completeness theorem (1929) showed that validity coincides with provability, the problem became: can a machine determine whether a given formula has a proof? Hilbert assumed the answer was yes --- another expression of his ``no \emph{ignorabimus}.'' What is striking in retrospect is that the problem could not even be \emph{precisely stated}, because the word ``procedure'' (\emph{Verfahren}) was undefined. What is a procedure? A recipe? A human habit? The word ``algorithm'' had no mathematical content. The Entscheidungsproblem thus turned the definition of computation from a philosophical puzzle into a concrete technical requirement: before asking whether a decision procedure exists, one must say exactly what a procedure \emph{is}.

\subsection{Mathematics Before Computers}
\label{sec:computation:math_before_computers}

To appreciate the urgency of the Entscheidungsproblem, it helps to recall what mathematics looked like before machines. Calculation was a human activity, and ``computer'' was a job title --- a person who computes. Scientific tables (for navigation, astronomy, ballistics) were produced by armies of human computers, often women, following fixed rules with pencil and paper; their errors were a known and feared source of disaster. Babbage designed the Analytical Engine (1837), with its ``mill'' (arithmetic unit) and ``store'' (memory), to eliminate exactly that human error --- but the machine was never built. The point is not nostalgia: the \emph{human computer} is the conceptual raw material of Turing's definition. When Turing set out to define computation, he did not abstract from electronics (none existed); he abstracted from the human clerk, watching what such a clerk does --- read a symbol, write a symbol, shift attention, follow a finite rule --- and stripped it to its logical core. The computer-as-person is the bridge between mathematics by hand and the machine.

\subsection{Why This Question Changed History}
\label{sec:computation:why_changed_history}

The Entscheidungsproblem turned out to be the hinge on which the digital age turned --- because its answer was the modern computer. When Turing answered it in 1936, he did so by constructing a mathematical model of computation and then establishing two things. First, no machine of this kind can decide validity --- the problem is unsolvable (Chapter 2). Second --- and this is the part nobody expected --- there is a single \emph{universal} machine that can simulate every other machine; a problem is ``computable'' precisely when this one machine can handle it. The decision problem, which asked whether all of mathematics could be reduced to calculation, was answered by inventing the very notion of calculation that underlies every computer. Hilbert asked whether truth could be mechanized; Turing showed what a mechanism for truth would \emph{be}. That is why a question posed in 1928 about the foundations of logic became, within fifty years, the theoretical foundation of an industry. The rest of this chapter develops that breakthrough.

Key developments converging in the 1930s:

\begin{itemize}
\item \textbf{Kurt G\"odel (1931):} Incompleteness theorems. Introduced \emph{G\"odel numbering} --- encoding syntax as arithmetic. Showed that provability can be arithmetized. But ``computable'' remained informal.

\item \textbf{Alonzo Church (1932--1936):} Lambda calculus. Defined \emph{$\lambda$-definable functions} as a formal counterpart to ``effectively calculable.'' Proposed \emph{Church's Thesis}: $\lambda$-definable = effectively calculable.

\item \textbf{Stephen Kleene (1935):} Recursive functions. Built on Herbrand and G\"odel's work. Defined \emph{general recursive functions} --- another formal class.

\item \textbf{Alan Turing (1936):} Turing machines. Analyzed what a human computer actually \emph{does} --- writes symbols on paper, changes mental state, moves left/right. Abstracted this into a mathematical machine.
\end{itemize}

G\"odel's encoding of syntax into arithmetic (Chapter 3) was the crucial precedent: it showed that \emph{meta-mathematical} notions (provability, truth) could be brought inside arithmetic. Turing applied the same insight to \emph{computation itself} --- the machine becomes data.

\subsection*{The Pre-1936 Landscape}

Before 1936, several mathematicians had circled the problem:
\begin{itemize}
\item \textbf{Emil Post (1921, 1936):} Developed ``canonical systems'' and ``tag systems'' --- string rewriting systems. \cite{church1936}. His 1936 paper (submitted before Turing's but published after) independently formulated a similar model. Post suffered from bipolar disorder, which delayed his work; his 1936 paper begins poignantly: ``The writer has been tempted to date the present paper 1921, but the temptation has been resisted.''

\item \textbf{Kurt G\"odel (1934):} In his Princeton lectures, G\"odel credited Herbrand for the notion of general recursive functions but expressed skepticism about Church's thesis. G\"odel's caution was scientifically justified: he wanted a compelling argument that the proposed formalism could not be extended.

\item \textbf{Alonzo Church (1932--36):} The $\lambda$-calculus emerged from Church's work on the foundations of logic. The (\cite{kleenerosser1935}), Kleene-Rosser paradox (1935) showed the original system was inconsistent; the ``pure'' $\lambda$-calculus survived as a model of computation. Church was Turing's PhD advisor, yet remarkably, Turing's approach was entirely independent.

\item \textbf{Alan Turing (1935--36):} As a 23-year-old graduate student at Cambridge, Turing was influenced by M. H. A. Newman's lectures on Hilbert's problems. His insight came from \emph{analyzing the human computer} rather than from logic. Turing's paper was submitted to the London Mathematical Society on May 28, 1936, received on May 30, and published in 1937.
\end{itemize}

\subsection*{Biographical Sketches}

Turing was a British mathematician, logician, and pioneer of computer science. He earned a first-class degree in mathematics from King's College, Cambridge in 1934, and was elected a Fellow at age 22. His 1936 paper, written when he was 23--24, is arguably the single most important paper in the history of computer science. During World War II, Turing worked at Bletchley Park, leading efforts to break the German Enigma code. He designed the Bombe, an electromechanical device that automated cryptanalysis. After the war, he worked on the Automatic Computing Engine (ACE) at the National Physical Laboratory. He was convicted of gross indecency for his homosexuality in 1952, subjected to chemical castration, and died of cyanide poisoning in 1954 in what was ruled a suicide. He was posthumously pardoned in 2013.

Church was an American mathematician and logician who spent most of his career at Princeton University. He invented the $\lambda$-calculus in 1932 as part of a formal system for the foundations of mathematics. Church was Turing's PhD advisor; their student-advisor relationship is one of the most consequential in intellectual history --- both men independently defined computation, and Turing proved their formalisms equivalent. Church also formulated Church's theorem: there is no decision procedure for first-order logic.

G\"odel was an Austrian--American logician, mathematician, and philosopher. His incompleteness theorems (1931), published when he was 25, are among the most important results in mathematical logic. G\"odel's work on recursive functions, developed in his 1934 Princeton lectures, was foundational to computability theory. G\"odel was a perfectionist who published relatively little, but each publication was transformative. He suffered from mental health problems in later life and died of starvation in 1978.

Kleene was an American mathematician who earned his PhD under Church at Princeton (1934). He developed the theory of general recursive functions independently of and concurrently with G\"odel and Herbrand. Kleene also founded the theory of regular languages and finite automata (Kleene's theorem, 1956) and introduced the Kleene star operator. His 1952 book \emph{Introduction to Metamathematics} was the definitive textbook on recursion theory for decades.

Post was an American mathematician and logician. He earned his PhD from Columbia (1920) for a dissertation on the propositional calculus. Post developed his canonical systems in 1921 but was unable to publish due to mental health struggles. He independently formulated a model equivalent to Turing machines in 1936, though his paper appeared slightly after Turing's. Post's ``tag systems'' --- simple string rewriting systems --- were later shown to be universal. He also proved the completeness and consistency of the propositional calculus, independent of Wittgenstein's contemporaneous work.

\subsection*{Beyond the Mainstream}

The problem of defining computation was not confined to the logicians. Several other threads converged:

\begin{itemize}
\item \textbf{Jacques Herbrand (1931):} Herbrand, a brilliant French mathematician, proposed a definition of recursive functions in correspondence with G\"odel. He died at age 23 in a mountaineering accident before his work was published.
\item \textbf{Charles Babbage (1837):} Designed the Analytical Engine, the first general-purpose computer design, a century before Turing. Babbage's machine was mechanical and never fully built, but its architecture anticipated the stored-program concept.
\item \textbf{Ada Lovelace (1843):} Wrote the first algorithm intended for machine execution (Bernoulli numbers on the Analytical Engine). Lovelace recognized that the machine ``weaves algebraic patterns, not just numbers.'' She also glimpsed the limits: ``The Analytical Engine has no pretensions whatever to originate anything.''
\item \textbf{David Hilbert (1926):} In ``On the Infinite,'' Hilbert wrote about the ``finite standpoint'' (\emph{finite Einstellung}) --- the idea that metamathematical reasoning must be based on finitary methods. This directly inspired the search for a finitary definition of computation.
\end{itemize}

The question ``What is computation?'' did not arise in a vacuum. As the subsections above trace, it grew from two millennia of work on algorithmic problem-solving --- and it was the Entscheidungsproblem that forced the notion to be made precise.

%======================================================================
% SECTION 4: THE MENTAL LEAP
%======================================================================
\section{The Mental Leap: The Creative Insight That Changed Everything}
\label{sec:computation:leap}

Turing's 1936 paper ``On Computable Numbers, with an Application to the Entscheidungsproblem'' contains one of the most extraordinary intellectual leaps in history. He did not \emph{propose} a definition of computation. He \emph{analyzed} what a human computer does and proved that his machine captures it.

\textbf{The Insight:} A human computer performing a calculation is a finite-state machine reading and writing symbols on an unbounded tape. The ``mental state'' of the computer is finite. The symbols are from a finite alphabet. The operations are local: read a symbol, change state, write a symbol, move left or right. \emph{That is all computation is.}There is no ``understanding'' in the machine. There is only state transition. The meaning lives entirely in the \emph{interpretation} we give to the input and output. The machine itself is pure syntax.

This leap had three audacious components:

\begin{enumerate}
\item \textbf{Reductionism:} The infinite variety of mathematical calculations reduces to a single, simple machine architecture.

\item \textbf{Universality:} There exists a \emph{Universal Turing Machine} $U$ that, given the description of any Turing machine $M$ and input $x$, simulates $M(x)$. One machine to rule them all.

\item \textbf{Equivalence as Evidence:} Turing didn't just define his machines. He proved they coincide with Church's $\lambda$-calculus and Kleene's recursive functions. Three independent formalisms, born from completely different intuitions, converged on the \emph{same class} of functions. This convergence is the strongest evidence that the class is \emph{natural} --- not an artifact of any one formalism.
\end{enumerate}

Before Turing, ``universal'' meant ``applies to everything in a domain.'' After Turing, ``universal'' meant ``a single machine that can simulate any other machine in the domain.'' The Universal Turing Machine is a \emph{programmable computer} --- and it was described before any physical computer existed. The abstraction predated the reality.

\subsection*{The Human Computer Analysis}

Turing's analysis in Section 9 of his 1936 paper is worth quoting at length. He imagined a human computer with:
\begin{itemize}
\item A two-dimensional paper (which can be reduced to a one-dimensional tape)
\item A finite set of ``states of mind''
\item The ability to observe a bounded number of squares at once
\item Operations: change observed symbol, change state, move left/right
\end{itemize}

He argued that any ``effective procedure'' a human could follow must be reducible to these operations. This \emph{empirical} approach --- deriving the model from observation of human behavior --- distinguished Turing from Church and G\"odel, who worked from purely mathematical intuitions.

\subsubsection*{Turing's Argument in Detail}

Turing's argument had seven steps:

\begin{enumerate}
\item \textbf{Discrete states:} The human computer's ``state of mind'' changes only when a symbol is written or observed. Since the number of symbols the computer can distinguish is finite (the human eye cannot distinguish infinitely many symbols), the number of states of mind is also finite.

\item \textbf{Bounded observation:} The human computer can only observe a bounded number of squares at once. Turing argued that the computer's field of vision is bounded by physiological limits --- we cannot scan arbitrarily many squares simultaneously.

\item \textbf{Locality of action:} The computer's action depends only on the currently observed symbols and the current state of mind. This follows from the nature of rule-following: if the computer is following a finite set of instructions, the next action must be determined by what is currently observable.

\item \textbf{Tape reduction:} A two-dimensional sheet of paper can be reduced to a one-dimensional tape by a space-filling curve or by numbering the cells sequentially. The human computer can treat the tape as a linear sequence.

\item \textbf{Symbolic reduction:} The computer's paper contains only symbols from a finite alphabet. Even if the computer writes in natural language, each meaningful expression can be decomposed into a finite set of atomic symbols.

\item \textbf{Determinacy:} For an effective procedure, the next action must be uniquely determined by the current state of observation. If multiple actions are possible, the procedure is not fully specified.

\item \textbf{Simulation:} The combination of these constraints yields a machine whose behavior is exactly that of a Turing machine. Therefore, any effective procedure can be executed by a Turing machine.
\end{enumerate}

Turing's analysis inverted the problem. Instead of asking ``What can a machine compute?'' he asked ``What does a human do when computing?'' By answering the second question, he answered the first. The human computer became the template for the machine computer.

\subsubsection*{A Worked Example: Adding Two Numbers by Hand}

Consider what happens when a human adds 472 and 389 using paper and pencil:

\begin{enumerate}
\item Write the numbers aligned to the right:
\[
\begin{array}{r}
  472 \\
+ 389 \\
\hline
\end{array}
\]
\item Start from the rightmost column: $2 + 9 = 11$. Write $1$, carry $1$.
\item Move left: $7 + 8 + 1 = 16$. Write $6$, carry $1$.
\item Move left: $4 + 3 + 1 = 8$. Write $8$.
\end{enumerate}

Turing's analysis decomposes this into TM operations:
\begin{itemize}
\item The paper is the tape, the columns are cells.
\item The ``write, move left'' sequence is a TM transition.
\item The carries are remembered in the finite state ($q_{carry}$ and $q_{nocarry}$).
\item The algorithm is a finite table of state-symbol rules.
\end{itemize}

Every arithmetic procedure, no matter how sophisticated, decomposes similarly. This decomposition is what Turing proved always possible.

\subsection*{Universality: The Machine That Contains All Machines}

The leap from ``a Turing machine can compute any computable function'' to ``there is a single Turing machine that can simulate any other'' is not obvious. It requires that the machine's description be \emph{data} --- that we can encode the state table of $M$ as a string on $U$'s tape, and that $U$ can interpret this string as instructions.

\textbf{The Second Insight:} If a Turing machine can compute any computable function, and if the function of ``simulating another Turing machine'' is itself computable, then there must be a Turing machine that does it. This is the universal machine. The machine that can read its own program --- stored in the same memory as its data.
This insight --- the stored program --- is the foundation of every modern computer. It means that hardware and software are separable: the same physical machine can become a word processor, a video game, or a scientific simulator, simply by loading different data into its memory.

%======================================================================
% SECTION 5: THE MATHEMATICS
%======================================================================
\section{The Mathematics: Rigorous Development of the Theory}
\label{sec:computation:math}

\subsection{Definitions}
\label{sec:computation:definitions}

\begin{definition}[Turing Machine]
A \emph{Turing machine} is a 7-tuple $M = (Q, \Sigma, \Gamma, \delta, q_0, q_{accept}, q_{reject})$ where:
\begin{itemize}
\item $Q$ is a finite set of states
\item $\Sigma$ is the input alphabet (not containing the blank symbol $\sqcup$)
\item $\Gamma$ is the tape alphabet, $\Sigma \cup \{\sqcup\} \subseteq \Gamma$
\item $\delta: Q \times \Gamma \to Q \times \Gamma \times \{L, R\}$ is the transition function
\item $q_0 \in Q$ is the start state
\item $q_{accept}, q_{reject} \in Q$ are halting states ($q_{accept} \neq q_{reject}$)
\end{itemize}
A \emph{configuration} is a string $uqv$ where $u,v \in \Gamma^*$, $q \in Q$, representing tape contents with the head at the first symbol of $v$ in state $q$.
\end{definition}

\begin{definition}[Computable Function]
A function $f: \Sigma^* \to \Sigma^*$ is \emph{Turing-computable} if there exists a Turing machine $M$ such that for every input $w \in \Sigma^*$, $M$ halts on $w$ with $f(w)$ written on its tape.
\end{definition}

\begin{definition}[Universal Turing Machine]
A \emph{Universal Turing Machine} $U$ is a Turing machine such that for every Turing machine $M$ and input $x$, $U(\langle M \rangle, x) = M(x)$, where $\langle M \rangle$ is an encoding of $M$ as a string.
\end{definition}

\subsection{The Church-Turing Thesis}
\label{sec:computation:ctt}

The \emph{Church-Turing Thesis} (CTT) is not a theorem --- it cannot be proved because ``effectively calculable'' is an informal notion. It is a \emph{hypothesis} about the relationship between the formal and the intuitive:

\begin{quote}
\textbf{Church-Turing Thesis:} A function on strings is effectively calculable (in the intuitive sense) if and only if it is computable by a Turing machine (equivalently, $\lambda$-definable, or general recursive).
\end{quote}

\begin{historicalnote}[The Thesis That Became a Definition]
G\"odel initially rejected Church's thesis (1934), calling it ``thoroughly unsatisfactory.'' But after seeing Turing's analysis (1936), G\"odel wrote: ``Turing's work gives a \emph{definition} of the general notion of formal system\ldots This concept is of great importance for the foundations of mathematics.'' G\"odel's endorsement carried immense weight.
\end{historicalnote}

The CTT exists in several forms:
\begin{itemize}
\item \textbf{Church's Thesis (1936):} $\lambda$-definable = effectively calculable.
\item \textbf{Turing's Thesis (1936):} Turing-computable = effectively calculable.
\item \textbf{Post's Thesis (1936):} Canonical systems = effectively calculable.
\item \textbf{Physical CTT (Deutsch 1985):} Every finitely realizable physical system can be perfectly simulated by a universal Turing machine.
\item \textbf{Extended CTT:} Any ``reasonable'' model of computation can be simulated by a probabilistic TM with at most polynomial overhead.
\end{itemize}

\begin{example}[A Function That Is Not Effectively Calculable]
Consider the function $f: \N \to \{0,1\}$ defined by:
\[
f(n) = \begin{cases}
1 & \text{if the $n$th Turing machine halts on input $n$}\\
0 & \text{otherwise}
\end{cases}
\]
This is the halting function. We will prove in Chapter 2 that it is not Turing-computable. By the CTT, it is not effectively calculable by any intuitive procedure. This is the first example of a well-defined mathematical function that no human or machine can compute.
\end{example}

\subsection{Equivalence of Formalisms}
\label{sec:computation:equivalence}

\begin{theorem}[Equivalence of Computation Models]
The following classes of functions are identical:
\begin{enumerate}
\item Turing-computable functions
\item $\lambda$-definable functions (Church)
\item General recursive functions (Kleene/G\"odel/Herbrand)
\item Partial recursive functions
\item Functions computable by register machines (Shepherdson-Sturgis)
\item Functions computable by Markov algorithms
\item Functions computable by Post canonical systems
\end{enumerate}
\end{theorem}

\begin{proof}[Proof Sketch]
Circular implications:
\begin{itemize}
\item TM $\to$ $\lambda$-calculus: Encode TM configurations as $\lambda$-terms; transitions as $\beta$-reductions.
\item $\lambda$-calculus $\to$ Recursive: G\"odel numbering of $\lambda$-terms; $\beta$-reduction as primitive recursive operation.
\item Recursive $\to$ TM: Simulate primitive recursion and minimization on a TM tape.
\item TM $\to$ Register Machine: Tape $\leftrightarrow$ registers; state $\leftrightarrow$ program counter.
\item TM $\to$ Markov Algorithm: Tape symbols as alphabet; transitions as rewrite rules.
\item TM $\to$ Post System: Axioms as initial configuration; productions as transition rules.
\end{itemize}
Each simulation is effective and compositional.
\end{proof}

\subsubsection*{Detailed: TM $\leftrightarrow$ $\lambda$-Calculus}

The simulation of a Turing machine in the $\lambda$-calculus uses a representation of the tape as a pair of lists (left and right of head), encoded as $\lambda$-terms. The state is a $\lambda$-term that, given the current symbol and tape pair, returns the new state, symbol, and direction. The $\beta$-reduction of the $\lambda$-term corresponds to one TM step.

More concretely, we encode:
\begin{itemize}
\item Tape cells as Church numerals: blank is $0$, symbol $a_i$ is $i+1$.
\item The tape left of the head as a list $L = [l_k, l_{k-1}, \ldots, l_1]$, and right as $R = [r_1, r_2, \ldots]$, with the head at $r_1$.
\item A configuration $(q, L, r_1, R)$ as a $\lambda$-term $C = \lambda f. f\; q\; L\; r_1\; R$.
\item The transition function $\delta(q, r_1) = (q', s, D)$ as a rewriting rule on configurations.
\item For $D = L$, the new configuration is $(q', \text{tail}(L), \text{head}(L), s::\text{tail}(R))$.
\item For $D = R$, the new configuration is $(q', s::L, \text{head}(R), \text{tail}(R))$.
\end{itemize}

This construction was carried out by Turing in his 1937 PhD thesis under Church and later refined by others.

\subsubsection*{Detailed: $\lambda$-Calculus $\leftrightarrow$ Recursive Functions}

Kleene (1936) showed that every $\lambda$-definable function is general recursive by G\"odel-numbering $\lambda$-terms and showing that $\beta$-reduction is primitive recursive. The converse --- that every general recursive function is $\lambda$-definable --- uses the fact that the basic functions (zero, successor, projection) and the operations (composition, primitive recursion, minimization) can be encoded in the $\lambda$-calculus. The $Y$-combinator handles recursion.

\subsubsection*{Detailed: Register Machine Equivalence}

A register machine (Shepherdson-Sturgis 1963) has a finite set of registers $R_1, R_2, \dots$ holding natural numbers, and a program of instructions:
\begin{itemize}
\item $\text{INC}(i)$: $R_i := R_i + 1$
\item $\text{DECJZ}(i, L)$: if $R_i > 0$ then $R_i := R_i - 1$; else jump to label $L$
\end{itemize}

\textbf{TM $\to$ Register Machine:} Given a TM $M = (Q, \Sigma, \Gamma, \delta, q_0, q_{accept}, q_{reject})$, we simulate it on a register machine with 4 registers:
\begin{itemize}
\item $R_1$: Encodes the left portion of the tape (using G\"odel numbering).
\item $R_2$: Encodes the right portion of the tape.
\item $R_3$: Stores the current state $q \in Q$ as a number.
\item $R_4$: Scratch register for arithmetic.
\end{itemize}

The tape encoding works as follows. Let $\Gamma = \{\gamma_1, \ldots, \gamma_k\}$ with $\gamma_1 = \sqcup$ (blank). Assign each symbol a number $1,\ldots,k$. Then the tape content $s_1 s_2 \cdots s_m$ is encoded as the number:
\[
\text{encode}(s_1 s_2 \cdots s_m) = \sum_{j=1}^m \text{code}(s_j) \cdot (k+1)^{j-1}
\]
This is a base-$(k+1)$ representation of the tape. Reading the leftmost symbol (the one under the head) corresponds to computing $\text{encode}(R_2) \bmod (k+1)$. Writing a symbol corresponds to adjusting this value. Moving left pops from $R_1$ and pushes to $R_2$; moving right does the reverse.

\textbf{Register Machine $\to$ TM:} Given a register machine with $r$ registers, we simulate it on a TM with $r$ tapes (one per register), each storing the register value in binary. The program counter is tracked by the TM state. Increment and decrement are simple binary addition and subtraction. The jump-if-zero test requires checking whether the register tape contains all zeros.

\begin{example}[Register Machine for GCD]
The following register machine computes $\gcd(a,b)$:
\begin{verbatim}
L1: DECJZ(2, L4)   // if R2 = 0, done
    DECJZ(1, L3)   // if R1 = 0, swap
    // Compare: subtract registers
    INC(3)         // use R3 as temp
    DECJZ(1, L2)   // decrement both
    DECJZ(2, L2)
    JZ(0, L1)      // unconditional jump
L2: // R1 < R2 or R2 < R1 detected
    // (details omitted for brevity)
L3: // Swap R1 and R2
L4: // Halt, result in R2
\end{verbatim}
This demonstrates that register machines --- despite their minimal instruction set --- can compute any function computable by a Turing machine.
\end{example}

\subsection{The \texorpdfstring{$\lambda$}{λ}-Calculus: Syntax and Reduction}
\label{sec:computation:lambda}

The $\lambda$-calculus was developed by Alonzo Church in the early 1930s as part of a formal system for the foundations of mathematics. It is remarkably simple: there are only three kinds of terms, and three reduction rules.

\begin{definition}[$\lambda$-Terms]
The set of $\lambda$-terms $\Lambda$ is defined inductively:
\begin{enumerate}
\item \textbf{Variables:} Every variable $x, y, z, \ldots$ is a $\lambda$-term.
\item \textbf{Abstraction:} If $M$ is a $\lambda$-term and $x$ is a variable, then $(\lambda x. M)$ is a $\lambda$-term (representing a function of $x$).
\item \textbf{Application:} If $M$ and $N$ are $\lambda$-terms, then $(M\; N)$ is a $\lambda$-term (representing applying $M$ to $N$).
\end{enumerate}
\end{definition}

\begin{definition}[Reduction Rules]
\begin{enumerate}
\item \textbf{$\alpha$-conversion:} $\lambda x. M \to_\alpha \lambda y. M[y/x]$ where $y$ is not free in $M$. Renaming bound variables.
\item \textbf{$\beta$-reduction:} $(\lambda x. M) N \to_\beta M[N/x]$. Applying a function to an argument.
\item \textbf{$\eta$-reduction:} $\lambda x. (M x) \to_\eta M$ if $x$ is not free in $M$.
\end{enumerate}
A term is in \emph{$\beta$-normal form} if no $\beta$-reduction is possible. By the Church-Rosser theorem, if a term has a normal form, it is unique up to $\alpha$-conversion.
\end{definition}

\begin{example}[$\beta$-Reduction Step by Step]
Consider $(\lambda x. x\; x)(\lambda y. y)$:
\begin{align*}
(\lambda x. x\; x)(\lambda y. y)
&\to_\beta (x\; x)[\lambda y. y / x] \\
&= (\lambda y. y)(\lambda y. y) \\
&\to_\beta y[\lambda y. y / y] \\
&= \lambda y. y
\end{align*}
Each step substitutes the argument into the function body, eliminating one $\beta$-redex.
\end{example}

\subsubsection*{Church Numerals}

Church numerals encode natural numbers as functions that apply a function $n$ times:

\begin{definition}[Church Numerals]
\begin{align*}
\bar{0} &= \lambda f.\lambda x. x \\
\bar{1} &= \lambda f.\lambda x. f\; x \\
\bar{2} &= \lambda f.\lambda x. f\; (f\; x) \\
\bar{n} &= \lambda f.\lambda x. f^n\; x \quad \text{(apply $f$ $n$ times to $x$)}
\end{align*}
\end{definition}

\begin{example}[Church Numeral Arithmetic]
\textbf{Successor:} $\text{SUCC} = \lambda n.\lambda f.\lambda x. f\;(n\; f\; x)$

Applying SUCC to $\bar{2}$:
\begin{align*}
\text{SUCC}\; \bar{2}
&= (\lambda n.\lambda f.\lambda x. f\;(n\; f\; x))(\lambda f.\lambda x. f(f\; x)) \\
&\to_\beta \lambda f.\lambda x. f((\lambda f.\lambda x. f(f\; x))\; f\; x) \\
&\to_\beta \lambda f.\lambda x. f(f(f\; x)) \\
&= \bar{3}
\end{align*}

\textbf{Addition:} $\text{PLUS} = \lambda m.\lambda n.\lambda f.\lambda x. m\; f\;(n\; f\; x)$

\textbf{Multiplication:} $\text{MULT} = \lambda m.\lambda n.\lambda f. m\;(n\; f)$

\textbf{Exponentiation:} $\text{EXP} = \lambda m.\lambda n. n\; m$

Let us verify PLUS on $\bar{2}$ and $\bar{3}$:
\begin{align*}
\text{PLUS}\; \bar{2}\; \bar{3}
&= (\lambda m.\lambda n.\lambda f.\lambda x. m\; f\;(n\; f\; x))\; \bar{2}\; \bar{3} \\
&\to_\beta^* \lambda f.\lambda x. \bar{2}\; f\;(\bar{3}\; f\; x) \\
&\to_\beta^* \lambda f.\lambda x. (\lambda f.\lambda x. f(f\; x))\; f\;(\bar{3}\; f\; x) \\
&\to_\beta^* \lambda f.\lambda x. f(f(\bar{3}\; f\; x)) \\
&\to_\beta^* \lambda f.\lambda x. f(f(f(f\; x))) \\
&= \bar{5}
\end{align*}
\end{example}

\subsubsection*{Booleans and Conditionals}

Church also encoded Boolean logic:
\begin{align*}
\text{TRUE} &= \lambda x.\lambda y. x \\
\text{FALSE} &= \lambda x.\lambda y. y \\
\text{IF} &= \lambda b.\lambda x.\lambda y. b\; x\; y
\end{align*}

Note that IF TRUE $x$ $y$ reduces to $x$, and IF FALSE $x$ $y$ reduces to $y$, because TRUE selects the first argument and FALSE selects the second.

\subsubsection*{The $Y$-Combinator: Recursion Without Names}

The $Y$-combinator is a fixed-point combinator: for any function $F$, $Y F$ is a fixed point of $F$, meaning $F(Y F) = Y F$. This enables recursion in the $\lambda$-calculus, which has no built-in recursion mechanism.

\begin{definition}[$Y$-Combinator (Curry's variant)]
\[
Y = \lambda f.(\lambda x. f(x\; x))(\lambda x. f(x\; x))
\]
\end{definition}

\begin{proof}[Derivation of the Fixed-Point Property]
\begin{align*}
Y F &= (\lambda f.(\lambda x. f(x\; x))(\lambda x. f(x\; x))) F \\
    &\to_\beta (\lambda x. F(x\; x))(\lambda x. F(x\; x)) \\
    &\to_\beta F((\lambda x. F(x\; x))(\lambda x. F(x\; x))) \\
    &= F(Y F)
\end{align*}
Thus $Y F = F(Y F)$, confirming the fixed-point property.
\end{proof}

\begin{example}[Factorial via the $Y$-Combinator]
Define:
\[
\text{FACT} = Y\; (\lambda f.\lambda n. \text{IF (ISZERO}\; n)\; \bar{1}\; (\text{MULT}\; n\; (f\; (\text{PRED}\; n))))
\]

Let us trace the reduction for $n = \bar{3}$:

\begin{align*}
\text{FACT}\; \bar{3}
&= Y\; F\; \bar{3} \\
&= F(Y F)\; \bar{3} \\
&\to_\beta (\lambda f.\lambda n. \text{IF (ISZERO}\; n)\; \bar{1}\; (\text{MULT}\; n\; (f\; (\text{PRED}\; n)))) (YF)\; \bar{3} \\
&\to_\beta \text{IF (ISZERO}\; \bar{3})\; \bar{1}\; (\text{MULT}\; \bar{3}\; (YF\; (\text{PRED}\; \bar{3}))) \\
&\to_\beta \text{IF FALSE}\; \bar{1}\; (\text{MULT}\; \bar{3}\; (YF\; (\bar{2}))) \\
&\to_\beta \text{MULT}\; \bar{3}\; (YF\; \bar{2})
\end{align*}

At this point, $YF\; \bar{2}$ recurses, computing $2! = 2$, and then MULT $\bar{3}\; \bar{2} = \bar{6}$.

The $Y$-combinator is remarkable because it achieves recursion purely through self-application, without any named recursive construct. This is why the $\lambda$-calculus is a universal model of computation: a single reduction rule ($\beta$-reduction) suffices to express all computable functions.
\end{example}

\subsection{Fixed-Point Combinators and Recursion}
\label{sec:computation:fixedpoint}

The $Y$-combinator is the most famous fixed-point combinator, but there are others:

\begin{itemize}
\item \textbf{Turing's $\Theta$ combinator:} $\Theta = (\lambda x.\lambda y. y(x x y))(\lambda x.\lambda y. y(x x y))$
\item \textbf{The $Z$ combinator (for call-by-value):} $Z = \lambda f.(\lambda x. f(\lambda y. x x y))(\lambda x. f(\lambda y. x x y))$
\end{itemize}

Each of these provides a way to define recursive functions in a language with no built-in recursion.

Fixed-point combinators are the theoretical foundation of recursion in programming languages. In denotational semantics, the meaning of a recursive program is defined as the least fixed point of a continuous function on a complete partial order. The $Y$-combinator is the syntactic counterpart of this semantic fixed point.

\subsection{Smallest Universal Turing Machines}
\label{sec:computation:smallest}

Once Turing established the existence of a universal machine, researchers asked: how simple can a universal machine be? This search yielded progressively smaller machines:

\begin{itemize}
\item \textbf{Shannon (1956):} Proved that two internal states suffice for universality \emph{provided the number of symbols is large}, and dually that two symbols suffice \emph{provided the number of states is large}. He did not exhibit a small (state, symbol) pair; in fact a (2,2) machine cannot be universal, since the halting problem is decidable for such small machines.
\item \textbf{Minsky (1962):} Constructed a 7-state, 4-symbol UTM using tag systems, the smallest known at the time. (His 1967 book also described a 4-state, 7-symbol machine on a circular tape.)
\item \textbf{Rogozhin (1996):} Found minimal UTMs including a 4-state, 6-symbol machine, a 5-state, 5-symbol machine, and a 2-state, 18-symbol machine, and proved these are minimal in their classes.
\item \textbf{Neary and Woods (2009):} Constructed \emph{weakly} universal machines (universal only under a periodic input encoding), including a 2-state, 4-symbol machine and a 2-state, 3-symbol machine built on tag systems. Wolfram's 2-state, 3-symbol machine was proven (strongly) universal by Alex Smith (2007).
\item \textbf{Current frontier:} The smallest known strongly universal machine has state--symbol product $|Q| \cdot |\Gamma| = 10$ (Wolfram's 2-state, 5-symbol machine, universal by Smith's argument); allowing weak universality, Neary and Woods achieve the same product with 2 states and 4 symbols.
\end{itemize}

\begin{example}[Minsky's 7-State 4-Symbol UTM]
Minsky's machine operates on a circular tape (a loop). It uses:
\begin{itemize}
\item States: $Q = \{A, B, C, D, E, F, G\}$ (7 states)
\item Symbols: $\Gamma = \{\sqcup, 0, 1, a\}$ (4 symbols)
\item A complex encoding of the simulated machine's state and tape onto a circular tape
\end{itemize}
The machine works by simulating a 2-counter machine (which is Turing-complete) using a single tape with a special encoding. Each counter is represented by the distance between markers on the tape. The 7 finite states are sufficient to manage increment, decrement, and zero-test operations on both counters.
\end{example}

The search for the smallest UTM is not a mathematical curiosity. It reveals the fundamental complexity of universal computation --- exactly how much structure is needed to achieve universality. This relates to questions in physics (can the universe simulate itself?), biology (how simple must a self-replicating system be?), and the philosophy of mind (how simple a system can exhibit universal behavior?).

\subsection{Cellular Automata and Rule 110}
\label{sec:computation:ca}

A cellular automaton is a discrete dynamical system consisting of a grid of cells, each in a finite state, updated synchronously according to a local rule. Stephen Wolfram, in the 1980s, systematically investigated one-dimensional cellular automata with two states and nearest-neighbor rules --- the ``elementary'' cellular automata.

\begin{definition}[Elementary Cellular Automaton]
An elementary cellular automaton consists of:
\begin{itemize}
\item An infinite one-dimensional array of cells, each in state $0$ or $1$.
\item A local rule $f: \{0,1\}^3 \to \{0,1\}$ that determines the next state of a cell based on its current state and those of its two neighbors.
\item Synchronous update: all cells update simultaneously each time step.
\end{itemize}
Since there are $2^3 = 8$ possible neighborhood patterns and each maps to $\{0,1\}$, there are $2^8 = 256$ elementary rules. Each rule is numbered by the binary string of its outputs.
\end{definition}

\begin{example}[Rule 110]
Rule 110 is the elementary cellular automaton with rule table:
\[
\begin{array}{c|c}
\text{Neighborhood} (111, 110, 101, 100, 011, 010, 001, 000) & \text{Output} (0,1,1,0,1,1,1,0)
\end{array}
\]
The binary string $01101110_2 = 110_{10}$ gives the rule its name.
\end{example}

\begin{theorem}[Cook 2004]
Rule 110 is Turing-complete. That is, there exists a universal Turing machine implemented as a pattern of Rule 110 cells.
\end{theorem}

\begin{proof}[Proof Sketch]
The proof simulates a cyclic tag system (a variant of Post's tag systems) using Rule 110. The key insight is that Rule 110 supports three types of propagating structures (``gliders'') that interact to perform computation:
\begin{enumerate}
\item \textbf{Glider A:} A pattern of period 7 that moves left.
\item \textbf{Glider B:} A pattern of period 2 that moves right.
\item \textbf{Eater:} A stable pattern that absorbs gliders.
\end{enumerate}
Collisions between gliders encode logical operations. By carefully arranging initial conditions, one can simulate the operations of a tag system, which is known to be Turing-complete. Cook's proof was controversial --- it was originally presented in 2000 but rigorous peer review took four years due to the complexity of verifying the construction.
\end{proof}

Rule 110 demonstrates that computation can emerge from simple, local interactions. This connects to modern work on reservoir computing, neuromorphic hardware, and the question of whether complex computation can arise spontaneously in physical systems. It also connects to Chapter 7 (Kolmogorov Complexity): Rule 110 shows that simple rules can generate computationally complex behavior.

\subsection{The Universal Machine Construction}
\label{sec:computation:universal}

\begin{theorem}[Existence of Universal Turing Machine]
There exists a Turing machine $U$ such that for any Turing machine $M$ and input $x$, $U$ on input $\langle M, x \rangle$ simulates $M$ on $x$.
\end{theorem}

\begin{proof}[Construction Sketch]
Encode $M = (Q, \Sigma, \Gamma, \delta, q_0, q_{accept}, q_{reject})$ as a string $\langle M \rangle$:
\begin{itemize}
\item States $\to$ numbers $1..|Q|$
\item Tape symbols $\to$ numbers
\item Transition function $\delta$ as a lookup table: $(q, a) \mapsto (q', a', D)$
\end{itemize}
$U$ uses three tapes (simulatable on one):
\begin{enumerate}
\item Tape 1: $\langle M \rangle$ (read-only)
\item Tape 2: Simulated tape of $M$
\item Tape 3: Current state of $M$ + scratch
\end{enumerate}
Loop: Read current state and scanned symbol from Tape 2, look up transition in Tape 1, update Tape 2 and Tape 3. Repeat until halting state.

To reduce to a single tape, interleave the three tapes using a marked-cell encoding. Each cell on the single tape contains three symbols (one from each tape), distinguished by markers. The slowdown is at most a constant factor because the three-tape machine makes $O(T)$ steps and each step simulates in $O(1)$ on the single-tape machine.
\end{proof}

\subsubsection*{Encoding Details}

A concrete encoding scheme for $\langle M \rangle$ might use:
\begin{itemize}
\item States encoded in unary: state $q_i$ as a string of $i$ $1$'s.
\item Symbols encoded in unary: symbol $a_j$ as a string of $j$ $1$'s.
\item A transition $(q_i, a_j) \to (q_{i'}, a_{j'}, D)$ encoded as: $1^{i}01^{j}01^{i'}01^{j'}0D$, where $D = L$ or $R$.
\item The entire transition table written sequentially, separated by $00$.
\end{itemize}

\subsubsection*{The Stored-Program Concept}

The universal machine $U$ embodies the \emph{stored-program concept}: the program ($\langle M \rangle$) and the data ($x$) are both stored in the same memory (the tape), and the machine treats them uniformly. This is the architecture of every modern computer --- von Neumann explicitly credited Turing's universal machine as the inspiration for the EDVAC design (1945).

\subsubsection*{Universal Machine Variants}

\begin{itemize}
\item \textbf{Two-tape UTM:} Shannon (1956) showed a 2-symbol, 2-state UTM exists (but with complex encoding).
\item \textbf{Small UTMs:} Minsky (1962), Rogozhin (1996) found minimal UTMs: 4-state 6-symbol, 5-state 5-symbol, etc. The smallest known is 2-state 18-symbol (Rogozhin 1996).
\item \textbf{Self-interpreter:} A UTM written in its own language --- the basis of meta-circular evaluators in Lisp and compiler bootstrapping.
\end{itemize}

\subsection{Partial Recursive Functions}
\label{sec:computation:recursive}

An alternative approach, developed by Kleene and G\"odel, defines computability through recursive function theory.

\begin{definition}[Partial Recursive Functions]
The class of \emph{partial recursive functions} is the smallest class containing:
\begin{enumerate}
\item \textbf{Zero:} $Z(x) = 0$
\item \textbf{Successor:} $S(x) = x + 1$
\item \textbf{Projection:} $P_i^n(x_1, \ldots, x_n) = x_i$
\end{enumerate}
and closed under:
\begin{enumerate}
\item \textbf{Composition:} $f(\vec{x}) = h(g_1(\vec{x}), \ldots, g_k(\vec{x}))$
\item \textbf{Primitive recursion:}
\[
\begin{cases}
f(0, \vec{x}) = g(\vec{x}) \\
f(n+1, \vec{x}) = h(n, f(n, \vec{x}), \vec{x})
\end{cases}
\]
\item \textbf{Minimization:} $f(\vec{x}) = \mu y[g(y, \vec{x}) = 0]$, where $\mu y$ means ``the least $y$ such that $g=0$,'' and $f$ is undefined if no such $y$ exists.
\end{enumerate}
\end{definition}

\begin{example}[Primitive Recursive Definitions]
\textbf{Addition:}
\[
\begin{cases}
\text{add}(0, y) = y \\
\text{add}(n+1, y) = S(\text{add}(n, y))
\end{cases}
\]

\textbf{Factorial:}
\[
\begin{cases}
\text{fact}(0) = 1 \\
\text{fact}(n+1) = \text{mult}(n+1, \text{fact}(n))
\end{cases}
\]

\textbf{Ackermann Function:}
\[
\begin{cases}
A(0, n) = n + 1 \\
A(m, 0) = A(m-1, 1) \\
A(m, n) = A(m-1, A(m, n-1))
\end{cases}
\]
The Ackermann function is recursive but not primitive recursive --- it grows faster than any primitive recursive function. This shows that minimization (the $\mu$ operator) genuinely increases the expressiveness of the formalism.
\end{example}

The recursive functions framework will recur throughout this book: in Chapter 2 (the halting problem as a non-recursive set), in Chapter 4 (Rice's theorem as a property of index sets), and in Chapter 7 (Kolmogorov complexity, where the set of random strings is not recursively enumerable).

\subsection{Post Canonical Systems}
\label{sec:computation:post}

Emil Post's approach to defining computation was based on string rewriting. A Post canonical system is defined by:
\begin{itemize}
\item A finite alphabet $A$.
\item A finite set of axioms (initial strings).
\item A finite set of production rules of the form:
\[
\frac{g_{11}\alpha_1 g_{12}\alpha_2 \cdots g_{1n}\alpha_n g_{1,n+1}}{\;g_{21}\alpha_{\pi(1)} g_{22}\alpha_{\pi(2)} \cdots g_{2m}\alpha_{\pi(m)} g_{2,m+1}\;}
\]
where each $g_{ij}$ is a fixed string (the ``gap markers'') and each $\alpha_i$ is a variable string that can be substituted.
\end{itemize}

Post proved that his canonical systems are equivalent to Turing machines, providing yet another independent characterization of the computable functions.

\subsubsection*{Tag Systems}

A particularly simple variant is Post's tag system:
\begin{itemize}
\item A finite alphabet, a deletion number $d$, and a finite set of production rules $x \to w_x$.
\item Starting from a string $S$, repeatedly:
\begin{enumerate}
\item Append $w_x$ to the end of $S$, where $x$ is the first symbol of $S$.
\item Delete the first $d$ symbols of $S$.
\end{enumerate}
\item Halt when $S$ becomes empty or has length less than $d$.
\end{itemize}

Despite their extreme simplicity, tag systems with $d=2$ are Turing-complete (Minsky 1961), and Rule 110's universality was proved via a simulation of cyclic tag systems (Cook 2004).

%======================================================================
% SECTION 6: CONSEQUENCES
%======================================================================
\section{Consequences: What Became Possible}
\label{sec:computation:consequences}

The definition of computation changed everything:

\begin{enumerate}
\item \textbf{The Stored-Program Computer:} Von Neumann, influenced by Turing's universal machine, designed the EDVAC architecture (1945). The distinction between ``program'' and ``data'' vanished --- both are bits in memory. This is the architecture of every modern computer, from smartphones to supercomputers.

\item \textbf{Undecidability Results:} Once computation is formalized, we can ask: \emph{Which problems are computable?} The answer: \emph{Almost none.} The halting problem (Chapter 2), the Entscheidungsproblem, Hilbert's 10th problem (Matiyasevich 1970), the word problem for groups (Novikov 1955, Boone 1959) --- all undecidable. Each undecidability proof uses the same core technique: encoding Turing machine computations into the problem domain.

\item \textbf{Complexity Theory:} If computation is a mathematical object, we can measure its \emph{cost} --- time, space, randomness, communication. This birthed computational complexity (Chapter 8). The P vs NP problem, the most important open problem in computer science, asks whether problems whose solutions are easy to check are necessarily easy to solve.

\item \textbf{Programming Languages:} $\lambda$-calculus became the theoretical foundation of functional programming (Lisp, ML, Haskell, Scala, Clojure). The Curry-Howard correspondence links $\lambda$-calculus to logic: \emph{programs are proofs}. This means that a type checker is a proof verifier, and writing a program is doing mathematics.

\item \textbf{Artificial Intelligence:} If thinking is computation, can machines think? Turing's 1950 paper ``Computing Machinery and Intelligence'' reframed philosophy of mind as computer science. The Turing Test became a operational definition of intelligence. Modern AI --- neural networks, large language models, reinforcement learning --- all operate on the premise that intelligence is a computational phenomenon.

\item \textbf{Compiler Construction:} The universal machine is the theoretical basis for interpreters, compilers, and JITs. A compiler is a function from source language to target language --- both are computable functions. The theory of compilation reduces to: implement a Turing machine that translates $\langle M \rangle$ in language $L_1$ to $\langle M \rangle$ in language $L_2$.

\item \textbf{Formal Verification:} Model checking, theorem proving, program synthesis --- all treat programs as mathematical objects whose properties can be algorithmically analyzed. The universal machine is now a \emph{verifier}: given a program and a specification, does the program satisfy the specification?

\item \textbf{Blockchain and Smart Contracts:} The Ethereum Virtual Machine is a Turing-complete (bounded) state machine. The concept of ``code is law'' --- that contractual terms can be executed by an immutable program --- only makes sense because computation is a precisely defined, deterministic process.

\item \textbf{Cloud Computing and Virtualization:} A virtual machine is a software implementation of a universal Turing machine running on a physical one. The hypervisor simulates the guest OS just as $U$ simulates $M$. Every ``hello world'' in a Docker container, every AWS EC2 instance, traces its lineage to Turing's universal machine.
\end{enumerate}

Each of these consequences depends on the foundational insight that computation can be formalized. Without a precise definition, none of them would be possible --- we could not prove undecidability, measure complexity, build compilers, or design stored-program computers. The abstraction of computation is what makes the entire field of computer science possible.

%======================================================================
% SECTION 7: PHILOSOPHICAL SIGNIFICANCE
%======================================================================
\section{Philosophical Significance: What It Taught Us About Knowledge, Computation, or Reality}
\label{sec:computation:philosophy}

\begin{enumerate}
\item \textbf{Syntax Suffices for Semantics:} The machine manipulates symbols without understanding. Yet it produces correct results. Meaning is not in the machine --- it is in the \emph{interpretation} we apply to its input/output. This is the \emph{Chinese Room} argument (Searle 1980) avant la lettre. Searle argued that a human following a rulebook for manipulating Chinese characters (without understanding Chinese) is like a computer running a program --- syntax without semantics. But Turing's insight is that syntax \emph{suffices} for the production of correct behavior, even if it does not constitute understanding. The Chinese Room does not refute the possibility of strong AI; it merely clarifies that understanding is not required for computation.

\item \textbf{The Limits of Formalism:} Hilbert wanted to reduce mathematics to mechanical calculation. Turing showed that even \emph{mechanical calculation} has inherent limits (the halting problem). The dream of a complete, decidable mathematics is impossible --- not because we haven't found the right axioms, but because \emph{computation itself} has boundaries. This is a deep fact about the universe: there are well-posed mathematical questions that cannot be answered by any finite mechanical procedure.

\item \textbf{Universality as a Physical Principle:} The Universal Turing Machine is not just a mathematical construct --- it describes a \emph{physical possibility}. Any physical system that can implement the basic operations (state, symbol, transition) can simulate any other. This is the \emph{Physical Church-Turing Thesis} (Deutsch 1985): every finitely realizable physical system can be perfectly simulated by a universal computing machine. If the Physical CTT holds, then the universe is a computer --- and the limits of computation are the limits of physics.

\item \textbf{Computation as a Lens on Nature:} Once computation is a precise concept, we can ask: Is the universe a computer? (Zuse 1969, Fredkin 1990, Wolfram 2002). Is DNA a program? Is evolution an algorithm? The concept of computation escaped computer science and became a framework for physics, biology, and philosophy. Each of these questions attempts to export the precision of Turing's definition into domains where it may or may not fit.

\item \textbf{The Recursive Nature of CS:} Computer science is unique: the tool (the computer) is also the object of study. The universal machine can simulate itself --- the field is \emph{self-hosting}. This recursion is why CS advances exponentially: better tools build better tools. The compiler for a new language can be written in that language. The operating system for a new machine can be compiled on the old machine. This bootstrapping is possible only because of universality.

\item \textbf{The Church-Turing Thesis Shapes AI Debates:} The CTT implies that any computable cognitive function can be implemented on a Turing machine. Therefore, if human cognition is computable, machines can (in principle) replicate it. The debate then shifts to whether cognition \emph{is} computable --- a question that Turing's definition makes precise but does not answer. This framing has shaped AI research from 1950 to the present.
\end{enumerate}

%======================================================================
% SECTION 8: MODERN LEGACY
%======================================================================
\section{Modern Legacy: Where This Idea Appears Today}
\label{sec:computation:legacy}

\begin{itemize}
\item \textbf{Programming Language Theory:} Type systems, operational semantics, compiler correctness --- all grounded in formal models of computation. The $\lambda$-calculus is the assembly language of PL theory. The simply typed $\lambda$-calculus (Church 1940) is the foundation of modern type systems; System F (Girard 1972) extends it to parametric polymorphism, the basis of generics in Java, C++, and Rust.

\item \textbf{Quantum Computation:} The quantum Turing machine (Deutsch 1985) and quantum circuit model extend Turing's framework. The \emph{Extended Church-Turing Thesis} (that any reasonable physical computation can be efficiently simulated by a probabilistic TM) is challenged by quantum computers --- but the \emph{original} CTT (what is computable at all) still holds. Shor's algorithm (1994) factors integers in polynomial time --- a task believed to be intractable for classical computers --- demonstrating that quantum computation changes efficiency, not computability.

\item \textbf{Verification and Synthesis:} Model checking, theorem proving, program synthesis --- all treat programs as mathematical objects whose properties can be algorithmically analyzed. The universal machine is now a \emph{verifier}. Tools like Coq, Isabelle, and Lean allow mathematicians to write formal proofs verified by machines --- realizing Turing's vision of computation as the arbiter of mathematical truth.

\item \textbf{Biological Computation:} DNA computing (Adleman 1994), membrane computing (P\u{a}un 2000), neural computation --- all analyzed through the lens of Turing-equivalent (or sub-Turing) models. Adleman's experiment solved a 7-city instance of the Hamiltonian path problem using DNA molecules, demonstrating that computation is not restricted to silicon.

\item \textbf{Blockchain and Smart Contracts:} The Ethereum Virtual Machine is a Turing-complete (bounded) state machine. ``Code is law'' only makes sense because computation is a precisely defined, deterministic process. Smart contracts executed on a blockchain are the closest practical realization of the Universal Turing Machine: a program that runs unchanged on a distributed network, producing deterministic outputs from deterministic inputs.

\item \textbf{Compiler Bootstrapping:} A compiler for language $L$ written in $L$ --- the self-hosting pattern enabled by universality. The first C compiler (Kernighan and Ritchie, 1973) was written in C. The GNU Compiler Collection (GCC) compiles C using a C compiler. Rust's bootstrap chain (OCaml $\to$ Rust) relies on the same principle. Bootstrapping is a practical consequence of the universal machine: once we have a minimal implementation, we can use it to build richer implementations.

\item \textbf{Meta-circular Evaluators:} ``eval'' in Lisp, the SECD machine (Landin 1964), PyPy (Python in Python) --- universal machines in their own languages. A meta-circular evaluator is a Universal Turing Machine written in the language it interprets. This is the ultimate expression of the stored-program concept: the interpreter is both program and data.

\item \textbf{Hypercomputation Models:} Infinite time Turing machines (Hamkins-Lewis 2000), oracle machines (Turing 1939), trial-and-error predicates (Putnam 1965) --- exploring the boundary of the CTT. These models ask: what if we relax the constraints of the Turing machine (finite time, finite memory at each step)? They provide a rich taxonomy of ``beyond Turing'' computation, though whether any physical system realizes them remains speculative.

\item \textbf{Computational Biology:} The genome is a digital program (string over $\{A, C, G, T\}$), the ribosome is a reader/executor, and evolution is a stochastic optimization algorithm. This analogy --- made precise by Turing's definition --- has transformed biology into an information science. The Human Genome Project, CRISPR gene editing, and computational drug design all depend on the view of DNA as executable code.
\end{itemize}

\begin{itemize}
\item Chapter 2 (Halting Problem): The first fundamental limit of computation.
\item Chapter 3 (G\"odel): The mathematical precedent for encoding syntax as data.
\item Chapter 8 (Complexity): From ``what is computable?'' to ``what is efficiently computable?''.
\item Chapter 11 (Zero Knowledge): Computation as a tool for proving knowledge without revealing it.
\end{itemize}

%======================================================================
% SECTION 9: LESSONS FOR FUTURE SCIENTISTS
%======================================================================
\section{Summary}
\label{sec:computation:summary}

This chapter addressed the foundational question of what constitutes a computation, situating the inquiry within the historical context of Hilbert's Entscheidungsproblem. The chapter established three equivalent formalisms for effective calculability: Turing machines, $\lambda$-calculus, and general recursive functions. The Church--Turing Thesis was formulated as the claim that all three formalisms capture the same class of computable functions. The Universal Turing Machine was constructed, demonstrating that a single machine can simulate any Turing machine given its description as input, thereby establishing the theoretical basis for stored-program architectures. The chapter also examined the minimality of universal systems, presenting results on small universal Turing machines (e.g., Rogozhin's 4-state, 6-symbol machine) and the Turing-completeness of simple cellular automata (Rule 110, proved by Cook 2004). The equivalence between Turing machines and $\lambda$-calculus was discussed, along with the s-m-n theorem and the recursion theorem as foundational tools for self-reference in computation. The chapter concluded by surveying alternative models of computation---register machines, counter machines, tag systems, and Post canonical systems---and demonstrating their equivalence to Turing machines, thereby reinforcing the robustness of the Church--Turing Thesis. Several open problems were noted: the Physical Church--Turing Thesis, which asks whether physical processes can compute beyond Turing-computable functions; the Busy Beaver function $\Sigma(n)$, whose growth rate exceeds that of any computable function; and the ongoing search for minimal universal Turing machines.

\section*{Key Takeaways}
\begin{itemize}
\item The Church--Turing Thesis posits that Turing machines, $\lambda$-calculus, and general recursive functions define the same class of effectively calculable functions.
\item The Universal Turing Machine demonstrates that a single device can simulate any computation, providing the theoretical basis for general-purpose computers.
\item Multiple independent formalisms (Turing machines, $\lambda$-calculus, register machines, cellular automata) converge on the same class of computable functions, lending empirical support to the Thesis.
\item The s-m-n theorem and the recursion theorem provide the formal machinery for self-reference and program transformation in computability theory.
\item Open problems include the Physical Church--Turing Thesis, the exact growth rate of the Busy Beaver function, and the classification of minimal universal Turing machines.
\end{itemize}

%======================================================================
% EXERCISES
%======================================================================
% EXERCISES
%======================================================================
\section{Exercises}
\label{sec:computation:exercises}

\subsection*{Basic}

\begin{exercise}[Basic]
Prove that a Turing machine with a doubly-infinite tape (infinite in both directions) is equivalent in power to a standard Turing machine.
\end{exercise}

\begin{exercise}[Basic]
Show that a Turing machine with $k$ tapes can be simulated by a single-tape Turing machine with at most a quadratic slowdown in time.
\end{exercise}

\begin{exercise}[Basic]
Prove that a Turing machine with a \emph{stay} option (head can stay in place instead of moving left or right) is equivalent to the standard model.
\end{exercise}

\begin{exercise}[Basic]
Show that the class of regular languages is strictly contained in the class of languages recognizable by Turing machines. Provide a specific example of a language that is Turing-recognizable but not regular.
\end{exercise}

\begin{exercise}[Basic]
Write the Church numeral for 4. Verify that PLUS $\bar{2}$ $\bar{2}$ reduces to $\bar{4}$ by showing the $\beta$-reduction steps.
\end{exercise}

\subsection*{Intermediate}

\begin{exercise}[Intermediate]
Define a \emph{2-counter machine} (two unbounded registers, instructions: increment, decrement-and-branch-if-zero). Prove it is Turing-complete by simulating a 2-stack pushdown automaton.
\end{exercise}

\begin{exercise}[Intermediate]
Implement the $Y$-combinator in the $\lambda$-calculus and use it to define the factorial function. Show the reduction steps for $Y\;\text{fact}\;3$.
\end{exercise}

\begin{exercise}[Intermediate]
Prove that the class of functions computable by a register machine with only \emph{increment} and \emph{decrement-and-branch-if-zero} (no zero-test) is strictly weaker than Turing-complete.
\end{exercise}

\begin{exercise}[Intermediate]
Prove that the class of primitive recursive functions is strictly contained in the class of general recursive functions. (Hint: Show that the Ackermann function is recursive but not primitive recursive.)
\end{exercise}

\begin{exercise}[Intermediate]
Give an explicit encoding of a Turing machine with 3 states and 2 symbols as an input to a Universal Turing Machine. Describe the transition table encoding in detail.
\end{exercise}

\begin{exercise}[Intermediate]
Show that the Boolean functions AND, OR, NOT can be encoded in the $\lambda$-calculus using Church booleans. Then implement a 2-to-1 multiplexer.
\end{exercise}

\subsection*{Challenge}

\begin{exercise}[Challenge]
Research the \emph{Busy Beaver} function $\Sigma(n)$: the maximum number of 1s an $n$-state Turing machine can write on an initially blank tape before halting. Prove $\Sigma(n)$ grows faster than any computable function. What does this imply about the limits of formal systems?
\end{exercise}

\begin{exercise}[Challenge]
Research the \emph{Physical Church-Turing Thesis} (Deutsch 1985). What physical principles (relativity, quantum mechanics, thermodynamics) might limit the power of physical computation? Can analog computers (Pour-El Richards 1989) compute non-recursive functions?
\end{exercise}

\begin{exercise}[Challenge]
Prove the Church-Rosser theorem for the $\lambda$-calculus: if a term $M$ reduces to $N_1$ and $N_2$, then there exists a term $P$ such that both $N_1$ and $N_2$ reduce to $P$. Why does this imply that $\beta$-normal forms are unique when they exist?
\end{exercise}

\begin{exercise}[Challenge]
Construct a Turing machine that computes the Ackermann function $A(3,n)$ for arbitrary $n$. Estimate the number of steps required for $n=3$.
\end{exercise}

\begin{exercise}[Challenge]
Explore the relationship between Rule 110 and cyclic tag systems. Describe how a cyclic tag system can be embedded in Rule 110. Why does the proof of Rule 110's Turing-completeness require careful analysis of glider interactions?
\end{exercise}

\subsection*{Computational}

\begin{exercise}[Computational]
Implement a Universal Turing Machine simulator in Python. Encode the 6-state, 2-symbol Busy Beaver champion and verify it runs for 119,112,334,170,342,540 steps before halting (or as far as you can simulate).
\end{exercise}

\begin{exercise}[Computational]
Implement a register machine simulator. Write a register machine program that computes the GCD of two numbers using the Euclidean algorithm.
\end{exercise}

\begin{exercise}[Computational]
Implement a $\lambda$-calculus interpreter with the $Y$-combinator. Define Church numerals and implement addition, multiplication, and exponentiation. Verify that your interpreter correctly reduces $Y\;\text{fact}\;3$ to $\bar{6}$.
\end{exercise}

\begin{exercise}[Computational]
Implement a simulator for elementary cellular automata (256 rules). Demonstrate that Rule 30 produces chaotic patterns, Rule 90 produces Sierpinski's triangle, and Rule 110 produces complex propagating structures.
\end{exercise}

\begin{exercise}[Computational]
Implement a Turing machine that recognizes the language $\{0^n 1^n \mid n \geq 0\}$. Run it on inputs of various lengths and verify correctness.
\end{exercise}

%======================================================================
% TIMELINE
%======================================================================
\section{Timeline}
\label{sec:computation:timeline}
\addcontentsline{toc}{section}{Timeline}

\begin{tabularx}{\linewidth}{|l|X|}
\hline
\textbf{Year} & \textbf{Event} \\ \hline
1837 & Babbage designs the Analytical Engine (first general-purpose computer) \\ \hline
1843 & Lovelace writes the first algorithm for the Analytical Engine \\ \hline
1900 & Hilbert's 23 problems; 10th problem (Diophantine equations) and Entscheidungsproblem posed \\ \hline
1921 & Post develops canonical systems (unpublished until 1936) \\ \hline
1928 & Hilbert and Ackermann pose the Entscheidungsproblem explicitly \\ \hline
1931 & G\"odel's Incompleteness Theorems; G\"odel numbering introduced \\ \hline
1932--34 & Church develops $\lambda$-calculus; proposes Church's Thesis \\ \hline
1934 & G\"odel rejects Church's Thesis as ``unsatisfactory'' \\ \hline
1935 & Kleene defines general recursive functions \\ \hline
1936 & Turing's ``On Computable Numbers''; Turing machines, Universal TM, halting problem \\ \hline
1936 & Post publishes canonical systems; independent formulation \\ \hline
1937 & Turing proves equivalence of TM and $\lambda$-calculus (PhD thesis under Church) \\ \hline
1943 & McCulloch-Pitts neural networks (finite automata) \\ \hline
1945 & Von Neumann's EDVAC report; stored-program architecture \\ \hline
1950 & Turing's ``Computing Machinery and Intelligence''; Turing Test \\ \hline
1956 & Shannon's 2-state, 2-symbol UTM; birth of artificial intelligence (Dartmouth workshop) \\ \hline
1962 & Minsky's 4-state, 7-symbol UTM \\ \hline
1963 & Shepherdson-Sturgis register machines \\ \hline
1965 & Hartmanis-Stearns: Time/Space complexity classes; computational complexity born \\ \hline
1985 & Deutsch: Quantum Turing Machine; Physical Church-Turing Thesis \\ \hline
1994 & Adleman: DNA computing (Hamiltonian path) \\ \hline
1996 & Rogozhin: minimal UTMs (4-state 6-symbol, 5-state 5-symbol) \\ \hline
2002 & Wolfram: \emph{A New Kind of Science}; Principle of Computational Equivalence \\ \hline
2004 & Cook: Rule 110 is Turing-complete \\ \hline
2009 & Woods and Neary: 2-state 4-symbol UTM (current record) \\ \hline
\end{tabularx}

%======================================================================
% CITATIONS
%======================================================================
\section{Citations}
\label{sec:computation:citations}

\begin{enumerate}
  \item \cite{turing1936}
  \item \cite{church1936}
  \item \cite{godel1931}
  \item \cite{kleene1936}
  \item \cite{post1936}
  \item \cite{vonneumann1945}
  \item \cite{deutsch1985}
  \item \cite{davis2000}
  \item \cite{hodges1983}
  \item \cite{minsky1967}
  \item \cite{rogozhin1996}
  \item \cite{cook2004}
  \item \cite{wolfram2002}
\end{enumerate}

%======================================================================
% INDEX ENTRIES
%======================================================================
\index{computation!definition of}
\index{Turing machine}
\index{Universal Turing Machine}
\index{Church-Turing Thesis}
\index{Entscheidungsproblem}
\index{Hilbert, David}
\index{G\"odel, Kurt}
\index{Church, Alonzo}
\index{Turing, Alan}
\index{Kleene, Stephen}
\index{Post, Emil}
\index{$\lambda$-calculus}
\index{recursive functions}
\index{computability}
\index{register machine}
\index{combinator calculus}
\index{Y-combinator}
\index{Church numeral}
\index{Physical Church-Turing Thesis}
\index{cellular automaton}
\index{Rule 110}
\index{Busy Beaver}
\index{tag system}
\index{stored-program concept}
\index{axiomatic method}
\index{Euclid}
\index{al-Khw\=arizm\=i}
\index{Leibniz, Gottfried Wilhelm}
\index{characteristica universalis}
\index{calculus ratiocinator}
\index{Boole, George}
\index{Frege, Gottlob}
\index{Peano, Giuseppe}
\index{Principia Mathematica}
\index{foundations crisis}
\index{intuitionism}
\index{Brouwer, L.~E.~J.}
\index{law of the excluded middle}
\index{Hilbert's program}
\index{human computer}
